\chapter{Recipes}\label{ch:recipes}
\begin{goals}
\textbf{In this chapter} you save an analysis as a recipe, run and edit it, write one from scratch in a form, read the YAML file behind it, and check every recipe in a project at once.
\end{goals}

You set up a grade–tonnage window once: the series, the density, the cutoffs, the units. Next month the model is new and the setup is the same. A \term{recipe} is that setup, saved under a name as a small text file in the project. Running it puts the analysis back exactly as you left it, on whatever data the project now holds.

This part of the book builds one thing, a rung at a time: a month-end you start with one click and can hand to a reviewer. The recipe is the first rung. A routine (\cref{ch:routines}) runs several in order, checks (\cref{ch:checks}) judge the results, and a run record (\cref{ch:records}) keeps the evidence. The grade–tonnage recipe you save in this chapter is still at work in the last of them.

\begin{figure}[htbp]\centering
\begin{tikzpicture}[font=\sffamily\small, node distance=5mm,
  rung/.style={draw=rule, fill=tint, rounded corners=1.5pt, minimum height=9mm, minimum width=27mm, align=center, text=ink},
  arr/.style={-{Stealth[length=2mm]}, rule, line width=0.6pt}]
\node[rung, draw=amber, line width=0.8pt] (a) {recipe\\[-1pt]{\footnotesize one analysis}};
\node[rung, right=of a] (b) {routine\\[-1pt]{\footnotesize many, in order}};
\node[rung, right=of b] (c) {checks\\[-1pt]{\footnotesize judged}};
\node[rung, right=of c] (d) {run record\\[-1pt]{\footnotesize the evidence}};
\draw[arr] (a) -- (b); \draw[arr] (b) -- (c); \draw[arr] (c) -- (d);
\end{tikzpicture}
\caption{The automation ladder. This chapter is the first rung; each later chapter adds one.}\label{fig:recipes-ladder}
\end{figure}

\section{Save as recipe}\label{sec:recipes-save}
A recipe is a file in the project's \file{recipes/} folder, so you need a project first: save one with \menu{File>Save project as…} (\cref{ch:projects}). Without a project the status line says so and nothing is saved.

The \ui{save as recipe…} button sits:
\begin{itemize}
\item beside \ui{Run} in the grade–tonnage, swath and stats windows;
\item in the Calculations window, where it saves the layer's whole set of ƒ columns;
\item at the foot of the \ui{select by solid} dialog;
\item in the decorations panel, where it saves a figure: the view, the section, the decorations, the background and the export resolution.
\end{itemize}
micro asks for a name, suggests one from the layer (``deposit grade-tonnage''), and writes \file{recipes/deposit-grade-tonnage.yaml}. The status line confirms the path.

\begin{tryit}
Open the demo project from the start screen and save it to a folder (\menu{File>Save project as…}, \cref{ch:projects}). The grade–tonnage window opens with the demo, on \file{deposit.csv}: \col{FE} as the grade in \%, \col{SG} as the density, so its left axis reads tonnes. (If you closed it, right-click \file{deposit.csv} and choose \menu{Analyse>Grade-tonnage…}, then set \code{SG} as the density and \ui{\%} as the unit.)
\begin{enumerate}
\item Press \ui{save as recipe…} and accept the name \code{deposit grade-tonnage}. The status line reads \emph{saved recipe “deposit grade-tonnage” → recipes/deposit-grade-tonnage.yaml}.
\item Close the window.
\item Open \menu{Tools>Recipes}. The recipe is listed as \ui{deposit grade-tonnage · grade-tonnage}.
\item Choose \ui{Run} from its submenu. The window comes back, configured, with the curve drawn.
\end{enumerate}
\end{tryit}

Saving again under a name that already exists updates that file in place. That is how you change a recipe: open it, adjust the window, save it under the same name.

\section{Running, editing, renaming}\label{sec:recipes-menu}
\menu{Tools>Recipes} lists every recipe in the project with its tool. Each one has a submenu:

\begin{tabularx}{\linewidth}{@{}lL@{}}\toprule
\ui{Run} & opens the tool configured and runs it. It reads \ui{Run…} when the recipe asks questions first (\cref{sec:recipes-inputs}). \\
\ui{Run fresh (reuse nothing)} & routines without \code{inputs} only: runs every step again (\cref{sec:records-fresh}). \\
\ui{Edit (open configured)} & opens the tool with the recipe's settings and does not run it. A routine, or a tool without a window, opens in the routine editor instead (\cref{sec:routines-editor}). \\
\ui{Rename…} & changes the recipe's name. The file keeps its old file name. \\
\ui{Delete} & removes the file from the project, after a confirmation. This can't be undone. \\
\bottomrule
\end{tabularx}

What \ui{Run} does depends on the tool. Grade–tonnage, swath, stats, scatter / Q-Q, the cross-tab and validation open their window with the result. A resource-table recipe places its table in the composer. A calc-columns recipe applies its columns to the layer and reports how many. A figure recipe renders the figure and puts back everything it touched. A by-solid recipe runs the selection or the flag.

The other tools have no window to open: the QA reports, export, composite, pairing, grade shell, join, reconcile and section series. \ui{Run} runs such a recipe as a one-step routine (\cref{ch:routines}). The run tracker opens (\cref{sec:routines-tracker}), the step writes its table, its file or its layer, and the run leaves a record in \file{runs/} like any routine's (\cref{ch:records}). \ui{Edit (open configured)} shows that one-step routine in the routine editor and runs nothing.

The command palette has the same list: press \key{⌘K} (\key{Ctrl}+\key{K} on Windows and Linux) and type \code{run recipe}.

\subsection{When the layer isn't there}
A recipe names its layers \emph{by name}: \file{deposit.csv}, not ``the third layer''. If that layer isn't loaded, micro offers the loaded layers that could stand in: \ui{“deposit grade-tonnage” wants deposit.csv — use deposit\_2026-11.parquet}. Pick one and the recipe runs on it. If nothing suitable is loaded, the status line says what the recipe wanted.

\begin{note}
A single recipe may ask this one question. A routine may not: it never prompts, so every layer it names must be loaded before it starts (\cref{sec:routines-preflight}).
\end{note}

\section{Which tools can be recipes}\label{sec:recipes-tools}
\begin{tabularx}{\linewidth}{@{}lL@{}}\toprule
\code{tool} & what the recipe holds \\ \midrule
\code{gt} · \code{swath} · \code{stats} · \code{xy} · \code{xtab} & an analysis window's setup (\cref{ch:gt,ch:swath,ch:stats,ch:resources}) \\
\code{validation} · \code{pairing} & the model against its samples (\cref{ch:validation}) \\
\code{resource} & a resource table (\cref{ch:resources}) \\
\code{modelqa} · \code{dhqa} · \code{meshqa} & a QA report (\cref{ch:qa}) \\
\code{export} & a layer, or the records a filter keeps, as CSV or Parquet \\
\code{composite} · \code{shell} · \code{join} · \code{reconcile} · \code{sections} & tools that make a layer, columns or a file \\
\code{calccols} & a layer's ƒ columns (\cref{ch:filter}) \\
\code{solid} & select, flag or ratio by a solid or surface (\cref{ch:geometry}) \\
\code{figure} & a figure export (\cref{ch:figures}) \\
\code{compo} & one composer page (\cref{ch:composer,ch:atlas}) \\
\code{routine} & a list of the above (\cref{ch:routines}) \\
\bottomrule
\end{tabularx}

The first six rows are the seventeen tools a routine runs without a window (\cref{sec:routines-tools}); the form of \cref{sec:recipes-form} offers the same seventeen. A file whose \code{tool} isn't in this table is skipped when the project opens, and the status line says which file.

\section{The file}\label{sec:recipes-yaml}
Open \file{recipes/deposit-grade-tonnage.yaml} in any text editor:

\begin{lstlisting}[language=yaml]
# micro recipe — grade-tonnage. Hand-edit freely; series
# reference layers by NAME.
name: "deposit grade-tonnage"
tool: "gt"
params:
  series:
    - layer: "deposit.csv"
      col: "FE"
      unit: "pct"
      weight: "SG"
      filter: ""
      color: "#5b7fa6"
  nCut: 40
  scope: "all"
  gradeUnit: "pct"
  densUnit: "t/m3"
\end{lstlisting}

Three keys make a recipe: a \code{name}, a \code{tool}, and the tool's settings under \code{params}. A fourth, \code{inputs}, is optional (\cref{sec:recipes-inputs}).

The settings are the ones the window shows. For grade–tonnage, each \term{series} is a layer, a grade column, the grade's declared unit, a density expression (\code{weight}) and an optional filter. \code{nCut} is the window's \ui{cutoffs} field: the number of automatic cutoff steps between the data's lowest and highest grade. Display settings (a series' \code{color}, the axis ranges) ride along in the file and don't affect the numbers. \Cref{app:tools} lists every tool's settings.

\subsection{Stating the cutoffs}\label{sec:recipes-cutoffs}
The window chooses its cutoffs from the data's range. That is right for looking and wrong for a table you mean to set beside last month's: when the range moves, every cutoff moves with it, and no row has a partner. So the tool also takes an explicit list, which only the file can give:

\begin{lstlisting}[language=yaml]
params:
  cutoffs:
    - 30
    - 40
    - 50
    - 55
    - 60
    - 65
  gradeUnit: "pct"
\end{lstlisting}

Add the list to the file by hand (or write the recipe in the form of \cref{sec:recipes-form}, which has a \code{cutoffs} field). With \code{cutoffs} present, \code{nCut} is ignored. A routine \emph{requires} the list (\cref{sec:routines-rules}).

The file is yours to edit. Comment it, copy it into another project, keep an office set in a shared folder. When you drop a recipe file into \file{recipes/} by hand, \menu{File>Scan folder for new data…} picks it up (\cref{sec:projects-inbox}).

When you write YAML by hand, write it the way micro does: block style (one item per line, lists with dashes), with strings in double quotes.

\begin{why}
A recipe is data, never code. The YAML micro reads is a strict subset with no tags and no executable constructs, so a recipe file from a colleague can't carry a script. Running it is as safe as opening a spreadsheet of numbers, and a reviewer can read all of what it does.
\end{why}

\subsection{Settings are checked}
A setting micro doesn't know is an error, and the error names the nearest real one:

\begin{lstlisting}
grade-tonnage: no setting “cutofs” — did you mean “cutoffs”?
\end{lstlisting}

A wrong type is refused the same way (\code{nCut: expects a whole number, got "many"}), and so is a column the layer doesn't have. Nothing is silently ignored and no default quietly takes a typo's place.

\section{New recipe… — the form}\label{sec:recipes-form}
You don't have to run an analysis to record it. \menu{Tools>Recipes>New recipe…} opens a form generated from the tool's own list of settings.

\shot[0.62\linewidth]{recipes-form}{The new-recipe form for grade–tonnage. Required settings carry an asterisk; hover a label to read what the setting means.}

\begin{enumerate}
\item Choose the \ui{tool}. The form rebuilds for it.
\item Type a \ui{name}.
\item Fill in the settings. A layer setting is a picker over the loaded layers. A column setting suggests the chosen layer's columns and also takes an expression. A list of numbers is typed with commas or spaces (\code{30, 40, 50}). A list such as \code{series} has a \ui{+ series} button that adds a row.
\item Press \ui{Save recipe}.
\end{enumerate}

The form checks what you entered against the tool's list of settings, the same list a routine uses. If something is wrong, the problems appear in red above the button and nothing is saved.

What only a routine demands doesn't stop the save. A grade–tonnage recipe without \code{cutoffs} is a good recipe for the window, so the form saves it, and the status line says what a routine would still need: \emph{recipe “deposit grade-tonnage” saved — to use it in a routine, add: cutoffs: a routine must state it (automatic cutoffs come from the data’s range and move with it)}. A routine step that uses the recipe adds the cutoffs under \code{with:} (\cref{sec:routines-rules}).

The form lists the seventeen tools of \cref{sec:recipes-tools}. Figures, composer pages, calc columns and by-solid recipes are saved from their own windows.

\section{Recipes that ask first}\label{sec:recipes-inputs}
A\index{recipe!asking first} recipe can declare \code{inputs}: a short list of questions. Running it then opens a small dialog, and the answers land in the settings wherever a \code{\$name} marker stands.

\begin{lstlisting}[language=yaml]
name: "GT — any model"
tool: "gt"
inputs:
  - name: "model"
    label: "Block model"
    type: "layer"
  - name: "grade"
    label: "Grade column"
    type: "column"
    of: "model"
  - name: "density"
    label: "Density column"
    type: "column"
    of: "model"
params:
  layer: "$model"
  gradeUnit: "pct"
  cutoffs:
    - 30
    - 40
    - 50
    - 60
  series:
    - layer: "$model"
      col: "$grade"
      unit: "pct"
      weight: "$density"
\end{lstlisting}

\shot[0.5\linewidth]{recipes-inputs}{The dialog generated from \code{inputs}. The column pickers follow the layer chosen above them.}

Choose \ui{Run…} and the dialog shows one field per input: a picker of loaded layers for \code{model}, and for \code{grade} and \code{density} a picker of that layer's columns, which refills when you change the layer. Press \ui{Run}. Your answers are remembered for the session, so the next \ui{Run…} starts from them.

There are five input types (text, number, choice, layer, column). They work the same way in a routine, where they matter more, so the full table is in \cref{sec:routines-inputs}.

\begin{tip}
\code{inputs} is written by hand in the file. When you later re-save the recipe from its window under the same name, micro keeps the inputs you wrote.
\end{tip}

\section{Check all recipes…}\label{sec:recipes-checkall}
\menu{Tools>Recipes>Check all recipes…} reads every file in \file{recipes/} and checks it the way it will run, without running anything. (The item is listed once the project has a recipe.)

\shot[0.7\linewidth]{recipes-checkall}{Check all recipes: one line per file, with each problem listed beneath it.}

\begin{itemize}
\item A recipe for one of the seventeen tools has its settings checked against the tool's list. Unknown settings, wrong types and missing required settings are listed under it.
\item A routine is expanded and put through its pre-flight (\cref{sec:routines-preflight}): missing recipes, bad settings, columns that don't exist, two steps writing the same report.
\item A recipe for a tool without a settings list (a figure, a composer page, calc columns, by solid) is marked \emph{not checked}.
\end{itemize}

Each line starts with ✓ or ✗. A routine written for layers that aren't loaded isn't wrong, so missing layers are a note, not a problem: \emph{some layers aren't loaded — their columns weren't checked}.

\subsection{Old routine files}
Routines written before settings moved under \code{with:} (\cref{sec:routines-with}) are flagged as \emph{in the old step form}. The window then offers \ui{Rewrite old steps into with:}, which moves each step's settings under \code{with:} and writes the file back. The step's own keywords stay where they are.

\begin{tip}
Run \ui{Check all recipes…} after you rename a column or replace a model. It finds every recipe the change broke, in one pass, before a month-end run does.
\end{tip}
